\chapter{Lactose Intolerance}
\index{Lactose Intolerance}
\label{sec:Lactose Intolerance}

\section{Overview}
Lactose intolerance is a digestive disorder caused by lactase deficiency leading to an inability to digest lactose, the primary sugar in dairy products, resulting in gastrointestinal symptoms after ingestion. It affects approximately 65\% of the global adult population, with wide variation by ethnicity, being most prevalent in East Asian, African, and Native American populations.
\section{Classification}

\begin{description}[font=\normalfont\bfseries,leftmargin=3em,labelwidth=2.5em]
  \item[Cause] Metabolic/Enzyme deficiency
  \item[Organ System] Gastrointestinal
  \item[Pathophysiology] Primary lactase deficiency -> undigested lactose reaches colon -> bacterial fermentation -> gas, bloating, and osmotic diarrhea
  \item[Duration] Chronic
  \item[Onset] Gradual (adulthood)
  \item[Mode of Transmission] Non-communicable (genetic)
  \item[Clinical Course] Chronic, stable
  \item[Severity] Mild to Moderate
  \item[Extent of Disease] Localized
  \item[Age Group] Adults (most common)
  \item[Epidemiology] Affects 65\% of adults worldwide; varies by ethnicity (5\% in Northern Europeans to >90\% in East Asians)
  \item[ICD-11 Code] DC75.0
\end{description}

\section{Causes}
Primary lactose intolerance results from the genetically programmed decline in lactase-phlorizin hydrolase enzyme activity in the small intestinal brush border that occurs after weaning in most of the world's population. Secondary lactose intolerance can develop transiently from injury to the small intestinal mucosa due to conditions such as gastroenteritis, celiac disease, or Crohn's disease. Congenital lactase deficiency is a rare autosomal recessive disorder presenting in neonates. This condition develops from a combination of genetic and environmental factors that disrupt normal body processes. When these normal processes are disrupted, it leads to the symptoms and signs of the condition. Several factors can work together to trigger or make the condition worse. Knowing the specific cause helps your doctor choose the most effective treatment.

\section{Risk Factors}

\begin{itemize}[noitemsep]
  \item East Asian, African, Native American, or Hispanic ancestry
  \item Increasing age (lactase activity naturally declines after age 2--5)
  \item Premature birth (transient lactase deficiency)
  \item Conditions damaging the small intestinal mucosa (celiac disease, Crohn's disease)
  \item Family history of lactose intolerance
\end{itemize}

\section{Signs and Symptoms}

\subsection{Early symptoms}
\begin{itemize}[noitemsep]
  \item Early manifestations of lactose intolerance may be subtle and nonspecific
\end{itemize}

\subsection{Common symptoms}
\begin{itemize}[noitemsep]
  \item \subsection{Common symptoms}
  \item \textbf{Common symptoms:}
  \item Abdominal bloating, flatulence, and cramping within 30 minutes to 2 hours after consuming lactose-containing foods
  \item Watery diarrhea due to osmotic pull of undigested lactose in the colon
  \item Borborygmi (audible rumbling stomach sounds)
  \item Nausea, occasionally with vomiting after large lactose loads
  \item Pain or discomfort in the affected area
  \item Difficulty performing daily activities
\end{itemize}

\subsection{Less common symptoms}
\begin{itemize}[noitemsep]
  \item Atypical or less common presentations of lactose intolerance may occur
\end{itemize}

\subsection{Emergency warning signs}
\begin{itemize}[noitemsep]
  \item Severe or worsening symptoms of lactose intolerance require immediate medical evaluation
\end{itemize}
\section{Progression and Stages}
Lactose intolerance in primary form develops gradually after early childhood as lactase enzyme activity naturally declines, becoming clinically noticeable typically in adolescence or early adulthood. The condition is stable and non-progressive, with symptom severity depending on the amount of lactose consumed and the individual's residual lactase activity level.

\section{Complications}

\begin{itemize}[noitemsep]
  \item Nutritional deficiencies from avoidance of calcium and vitamin D-rich dairy products
  \item Increased risk of osteoporosis and fractures if dairy avoidance is not compensated
  \item Unintentional weight loss in severe cases
  \item Reduced quality of life
  \item Emotional and psychological effects
\end{itemize}

\section{Diagnosis}

\begin{itemize}[noitemsep]
  \item Hydrogen breath test after a 25--50 gram oral lactose load (gold standard)
  \item Lactose tolerance test measuring serial blood glucose levels
  \item Lab tests to check how your organs are working
  \item Imaging tests
\end{itemize}

\section{Prevention}

\begin{itemize}[noitemsep]
  \item Gradual introduction of lactose-containing foods to identify individual tolerance threshold
  \item Choosing lactose-reduced or lactose-free dairy alternatives
  \item Go for regular check-ups so any problems are caught early
  \item Follow your doctor's screening recommendations
\end{itemize}

\section{Treatment Options}

\begin{itemize}[noitemsep]
  \item Dietary lactose restriction based on individual tolerance
  \item Oral lactase enzyme supplements taken with meals
  \item Calcium and vitamin D supplementation to prevent bone loss
  \item Lifestyle changes and self-care
  \item Surgery when needed
\end{itemize}

\section{Living with It}

\begin{itemize}[noitemsep]
  \item Follow your treatment plan as prescribed
  \item Keep up with regular appointments with your healthcare provider
  \item Adjust your daily habits as needed
  \item Reach out to family, friends, or support groups for help
  \item Keep learning about your condition and how to manage it
\end{itemize}

\section{Outlook}
Lactose intolerance is a chronic but easily managed condition that does not reduce life expectancy. Your outlook depends on the specific condition, how advanced it is when found, and how well it responds to treatment. Getting treated early generally leads to better results. Many people manage this condition effectively with the right treatment. Regular check-ups are important to monitor your condition and adjust treatment as needed.
